We evaluate recorded deployments on Astropy instances from the test split of \texttt{princeton-nlp/SWE-bench\_Verified}, using Qwen Code 0.19.4 to inspect repositories, edit files, execute tools, and produce patches. A shared two-task comparison covers three tree runs and two completed SGLang EAGLE task pairs. A separate ten-task segment, Cqc10, recorded on 24 August 2026, supplies the latest complete single-configuration deployment case. The task/evaluation worker is x86-64; model serving runs on the GB10. These are full agent attempts. Table~\ref{tab:swe-tasks} lists every task and saved evaluator outcome in Cqc10.

\paragraph{Ten-task deployment}
The served model is the port's \texttt{qwen3.8-27b-nvfp4-radixark} checkpoint with ModelOpt mixed quantization and bfloat16 execution dtype. The route uses Hydra27 with 32 physical verification rows, full-vocabulary draft logits, draft-logit reuse, and a patched FA2 GQA-pair split-K4 attention kernel. Prefix caching and full/piecewise CUDA graphs are enabled. The engine permits one active sequence, with maximum context 131,072 tokens and a 4,096-token scheduler budget. Source commit \texttt{e7af6b595a8b} and the loaded extension hash are retained in the artifact. This graph/cache-enabled stochastic deployment is distinct from the Cat10 eager qualification below; passing that qualification does not prove this route's full-model equivalence.

The proxy supplies temperature 0.6, top-$p$ 0.95, top-$k$ 20, minimum-$p$ 0, and presence penalty 1.0; the engine seed is zero. The per-response output cap is 24,000 tokens, with a separate 20,000-token compaction cap and a 9,000-second agent timeout. Proxy auto-continuation is disabled. The agent runs in the task's repository image on an internal network that permits the model proxy and blocks external egress; per-task records retain the verified network-policy fingerprint. Agent duration comes from the task runner's elapsed-time field: it includes the agent loop and tool work, and excludes server startup and subsequent evaluation. We do not infer a completed-task speedup from model-only counters.

\begin{table}[t]
\caption{All ten tasks in the completed Cqc10 Astropy case. IDs have prefix \texttt{astropy\_\_astropy-}. Times are agent elapsed minutes, rounded from saved runner metadata; evaluation is subsequent. This selected development segment has one attempt per listed task and no matched native arm.}
\label{tab:swe-tasks}
\centering\small
\begin{tabular}{@{}lrl@{}}
\toprule
Task suffix & Agent time (min) & Saved evaluation \\
\midrule
13977 & 24.59 & Failed tests \\
14096 & 22.99 & Resolved \\
14182 & 7.23 & Failed tests \\
14309 & 1.67 & Resolved \\
14365 & 6.78 & Failed tests \\
14369 & 27.70 & Resolved \\
14508 & 41.24 & Resolved \\
14539 & 18.57 & Resolved \\
14598 & 25.79 & Failed tests \\
14995 & 4.59 & Resolved \\
\bottomrule
\end{tabular}
\end{table}

\paragraph{Workload observations}
All ten attempts have evaluator reports: six resolved their tasks and four failed tests. Recorded agent durations sum to 181.15 minutes, with a range of 1.67--41.24 minutes. This is an observed task ledger, not a representative SWE-bench Verified success estimate. The sum is accumulated agent time, not elapsed campaign time including startup and evaluation.

Every task used tools and produced a nonempty patch. A source-bound replay of the recorded trace checks detects no degeneration in these ten runs and no malformed tool-call arguments. The degeneration heuristic jointly requires at most 200 visible characters, zero tool calls, and at least 2,000 thinking characters; none of the ten traces meets that conjunction. This screen can miss loops that continue narrating or using tools, so its zero detections establish neither universal degeneration freedom nor preserved task quality. The earlier campaign's detected degeneration remains an adverse observation in the lineage audit.

\paragraph{Pooled decode comparison}
For completed request $r$, let $n_r$ be its output-token count, $\ell_r$ its server-recorded end-to-end latency, and $f_r$ its time to first output. With $R$ completed requests and $N=\sum_r n_r$, we use
\begin{equation}
 D_{\mathrm{pool}} =
 \frac{\sum_{r=1}^{R}(n_r-1)}{\sum_{r=1}^{R}(\ell_r-f_r)}
 = \frac{N-R}{\sum_r\ell_r-\sum_r f_r}.
 \label{eq:pooled-decode}
\end{equation}
Counts and durations are pooled before division. The numerator uses the conventional $n_r-1$ output intervals; outputs arriving in one batch can have zero arrival gaps. The denominator excludes first-output latency and between-request agent/tool work. It is accumulated request time, which counts overlapping requests separately, not isolated GPU time or elapsed campaign time.

For Cqc10, $N=248{,}342$, $R=265$, $\sum_r\ell_r=10{,}667.6211$ seconds, and $\sum_r f_r=989.8239$ seconds yield 25.63 tokens/s. The 265 requests reconcile with 256 normal requests and nine successful internal compaction requests; all recorded completion reasons are stop. The ten pre/post metric brackets are contiguous. Generation-token and completed-request counters agree with their latency-histogram populations, with no counter reset or outstanding request at the selected boundaries.

The later SGLang EAGLE comparator, recorded on 24 August, serves the as-shipped Qwen3.8-27B NVFP4 checkpoint with FlashInfer, using speculative steps/top-$k$/draft-token settings of 3/1/4. It uses the same Qwen Code version, agent harness, requested sampling settings, and network policy, with a 24,000-token response cap, a 262,144-token context limit, and an 8,192-token prefill chunk. Its two completed Astropy tasks are 12907 (resolved; 7.10 agent minutes) and 13033 (failed tests; 26.96 agent minutes). Their counters give $N=47{,}854$, $R=45$, $\sum_r\ell_r=2{,}013.6344$ seconds, and $\sum_r f_r=235.9488$ seconds, yielding 26.89 tokens/s. Both brackets contain streaming requests only. All completed engine requests are included, including internal agent traffic.

An earlier SGLang EAGLE run, recorded on 16--17 August, uses the same 3/1/4 speculative settings and a 32,768-token response cap. Its completed records for the same two task IDs give $N=47{,}850$, $R=17$, $\sum_r\ell_r=1{,}590.9038$ seconds, and $\sum_r f_r=32.8508$ seconds, yielding \textbf{30.70 tokens/s}. Both brackets pass the same completed-request, idle-boundary, and streaming-only checks. The harness runs one agent task at a time; the earlier SGLang record notes occasional overlapping auxiliary requests, which remain included in the summed request-time denominator. Task 12907 resolved; task 13033 produced an empty patch and was marked failed without invoking the test harness. Its final response contains 32,768 output tokens of thinking only, reaching the configured cap. We retain this faster observation and its unsuccessful agent outcome alongside the later run.

\paragraph{Best recorded shared-task result}
We select task IDs 12907 and 13033 because both have completed SGLang records, then report all three recovered split-K tree runs on that pair. Sr12 and Cqc16 ran on 19 August; Cqc15 ran on 23 August. Their engagement receipts identify Hydra27 with 27 valid draft nodes, maximum draft depth 11, 32 physical verification rows, full-vocabulary drafting, and GQA-pair split-K4 on all sixteen full-attention layers, with no fallback. They share the loaded attention extension but differ in Python patch revisions. Each uses a 32,768-token response cap, matching the earlier SGLang cap and exceeding the later run's 24,000. For the Sr12/later-SGLang pair, all 41 Sr12 responses and all 45 later SGLang responses are at most 20,000 tokens, below both configured caps; no response uses the differing output allowance. All Sr12 requests terminate with stop.

The best shared-task tree observation is Sr12: $N=25{,}027$, $R=41$, $\sum_r\ell_r=989.2612$ seconds, and $\sum_r f_r=130.4422$ seconds give \textbf{29.09 tokens/s}, 8.18\% above the later SGLang rate and 5.23\% below the earlier rate under Eq.~\ref{eq:pooled-decode}. Later observations are 28.28 for Cqc16 and 27.27 for Cqc15; Table~\ref{tab:pooled-decode} retains them alongside the best result. All selected task brackets pass the same population and idle-boundary checks. Every row on the shared pair has one resolved task and one failed task. All six selected tree attempts used tools and produced nonempty patches, with no detected degeneration or malformed tool-call arguments under the recorded checks. Every selected tree request finished with stop, without a recorded response-cap hit.

\begin{table}[t]
\caption{Pooled decode rates for tree verification and SGLang EAGLE with Qwen3.8-27B NVFP4 on GB10, using Eq.~\ref{eq:pooled-decode}. The first five rows use the same two SWE-bench Verified Astropy tasks, 12907 and 13033; their decoder settings are specified in the text. Cqc10 covers a separate ten-task subset. Bold marks each system's best recorded shared-task rate.}
\label{tab:pooled-decode}
\centering\small
\begin{tabular}{@{}lrrr@{}}
\toprule
Deployment & Tasks & Requests & Tokens/s \\
\midrule
SGLang early & 2 & 17 & \textbf{30.70} \\
SGLang later & 2 & 45 & 26.89 \\
Tree Sr12 (Aug 19) & 2 & 41 & \textbf{29.09} \\
Tree Cqc16 (Aug 19) & 2 & 35 & 28.28 \\
Tree Cqc15 (Aug 23) & 2 & 53 & 27.27 \\
\midrule
Tree Cqc10 (Aug 24) & 10 & 265 & 25.63 \\
\bottomrule
\end{tabular}
\end{table}

The tree headline is its best recorded rate on the shared pair, not the latest run or Sr12's entire workload. Sr12 completed four tasks in total, two resolved and two failed; pooling all four gives 26.21 tokens/s. Cqc16's subsequent task 13236 exhibited degeneration, and Cqc15's subsequent task 13398 lacks a completed evaluation. Those adverse records remain in the artifact and outside this explicitly two-task comparison. The shared-task comparison evaluates the complete decoding systems, including their proposal topology, draft budget, and serving optimizations. These method-specific choices are part of the comparison. The records use the same requested sampling settings and agent harness. We retain both recovered SGLang observations and all three recovered tree observations on this shared pair; they do not establish each system's optimum over draft lengths and candidate budgets. Repeated runs with a predeclared tuning protocol would measure that performance envelope; a component ablation would separately isolate the contribution of tree verification. The artifact retains the original metric brackets, source-bound configuration and behavior records, and reproducible pooled reductions.

\paragraph{Selection and comparison limits}
Cqc10 contains the ten remaining tasks of an interrupted sixteen-task development campaign. Earlier segments changed source revisions and output caps and included degeneration, incomplete output, and budget-capped attempts. Source-bound audit records of those failures, retries, and missing evaluations remain in the artifact; the original records remain at their recorded source paths. We do not relabel the completed segment as the entire campaign, combine those segments into a single-configuration benchmark score, or drop unsuccessful attempts from a claimed benchmark denominator. The tasks were selected from one repository, and this deployment has one recorded boot without a matched ten-task native run. Separate targeted native probes change other serving settings and do not supply that comparator. Consequently, these records demonstrate an implemented verifier serving real coding-agent work, but establish no native-versus-tree task-speed or quality advantage. A current, matched task campaign is the remaining experiment for that claim.


\paragraph{Existing comparator experiments}
Earlier SGLang EAGLE trials and native MTP controls remain in the repository. The earlier native campaign predates the NVFP4 deployment; later NVFP4 native probes cover targeted tasks with different serving settings. DSpark measurements use synthetic requests, and a later DFlash2 check establishes boot and response handling rather than agent-task performance. These records are distinguished from the pooled SWE comparison above; their superseded or non-agent rates are not additional performance results for the current verifier.
